\subsection{The Exponential and Gamma Distributions}
\hrulefill

\subsubsection*{The Exponential Distribution}
A continuous random variable $X$ is said to have an Exponential distribution with parameter
$\lambda > 0$ if its probability density function is given by
\[f(x) = \begin{cases}
\lambda e^{-\lambda x} & x > 0 \\
0 & \text{otherwise}
\end{cases}.\]

\paragraph*{Notation:} $X \sim \text{Exp}(\lambda)$ means that $X$ is an Exponential random variable with parameter
$\lambda$.

\paragraph*{CDF:} The CDF of $X \sim \text{Exp}(\lambda)$ is
\[F(x) = P(X \leq x) = \begin{cases}
1 - e^{-\lambda x} & x \geq 0 \\
0 & x < 0
\end{cases}.\]

\paragraph*{Used for:}
\begin{itemize}
    \item The time until a radioactive particle decays.
    \item The time until a customer arrives at a service point.
    \item The time until a light bulb burns out.
    \item Life spans of living organisms.
\end{itemize}

\paragraph*{Survival Function:} The survival function is given by
\[S(t) = P(X > t) = 1 - F(t) = e^{-\lambda t}.\]
This is the probability that the event of interest has not occurred by time $t$.
This is done instead of the CDF for lifespans because people are little bitches about death, so now we have to 
learn more formulae.

\paragraph*{Mean and Variance:} If $X \sim \text{Exp}(\lambda)$, then
\[\mu = E(X) = \frac{1}{\lambda}, \quad \sigma^2 = Var(X) = \frac{1}{\lambda^2}.\]

\paragraph*{Note:} When integrating by parts, let $u$ be the polynomial term and $dv$ be the exponential term.
\paragraph*{Note:} $\lim_{b\to\infty} \frac{\text{Polynomial}}{e^{\lambda b}} = 0$ for any polynomial.

\paragraph*{Lack of Memory Property}
The Exponential distribution has the lack of memory property, which states that
\[P(X > s + t | X > s) = P(X > t) \forall s,t \geq 0.\]
This means that the probability of waiting an additional time $t$ does not depend on how much time $s$ has already elapsed.

\paragraph*{Example:}
You are a customer service representative. Let $X$ be the time in minutes that you spend with a customer.
Suppose $X \sim \text{Exp}(\lambda = 0.5)$.\\
\\
a. What is the probability that it takes the customer representitive more than 3 minutes with the next customer?
\begin{align*}
    P(X > 3) &= 1 - F(3) = e^{-0.5 \cdot 3} = e^{-1.5} \approx 0.2231.
\end{align*}\\

b. Suppose the rep has already been on the phone for 15 minutes. What is the probability it will take a total of 18 minutes 
or more with this customer?
\begin{align*}
    P(X > 18 | X > 15) &= \frac{P(X > 18 \cap X > 15)}{P(X > 15)} = \frac{P(X > 18)}{P(X > 15)} \\
    &= \frac{e^{-0.5 \cdot 18}}{e^{-0.5 \cdot 15}} = e^{-0.5 \cdot 3} = e^{-1.5} \approx 0.2231.
\end{align*}\\

c. What is the probability that it takes the customer rep less than 3 minutes with the next5 customers?
\begin{align*}
    &X_i \sim \text{Exp}(0.5) \quad i = 1,2,3,4,5& \\
    P(X_1 < 3 \cap X_2 < 3 \cap X_3 < 3 \cap X_4 < 3 \cap X_5 < 3) &= P(X_1 < 3)P(X_2 < 3)P(X_3 < 3)P(X_4 < 3)P(X_5 < 3) \\
    &= (F(3))^5 = (1 - e^{-0.5 \cdot 3})^5 \approx 0.6321^5 \approx 0.1013.
\end{align*}\\

d. What is the probability that it takes the rep a total of 15 minutes or less to help the next 5 customers?
\begin{align*}
    &X_i \sim \text{Exp}(0.5) \quad i = 1,2,3,4,5& \\
    T &= \sum_{i=1}^5 X_i \sim \Gamma(n=5, \lambda=0.5)\\
    P(T \leq 15) &= F_T(15) = 1 - e^{-0.5 \cdot 15} \sum_{k=0}^{4} \frac{(0.5 \cdot 15)^k}{k!} \\
    &= 1 - e^{-7.5} \left(1 + 7.5 + \frac{7.5^2}{2} + \frac{7.5^3}{6} + \frac{7.5^4}{24}\right) \approx 0.5596.
\end{align*}

\hrulefill

\subsubsection*{The Gamma Distribution}
\paragraph*{Definition} A continuous random variable with PDF
\[f(x) = \begin{cases}
\frac{\lambda (\lambda x)^{n-1} e^{-\lambda x}}{(n-1)!} & x > 0 \\
0 & \text{otherwise}
\end{cases}\]
is said to have a Gamma distribution with parameters $n$ and $\lambda$, where $n$ is a positive integer
and $\lambda > 0$.

\paragraph*{The Gamma Function}
The Gamma function is defined as
\[\Gamma(\alpha) = \int_0^\infty x^{\alpha - 1} e^{-x} \, dx \quad \alpha > 0.\]
For positive integers $n$, $\Gamma(n) = (n-1)!$.

\paragraph*{Properties of the Gamma Function:}
\begin{itemize}
    \item For $\alpha > 1, \Gamma(\alpha) = (\alpha - 1) \Gamma(\alpha - 1)$.
    \item For a positive integer $n, \Gamma(n) = (n-1)!$.
    \item $\Gamma\left(\frac{1}{2}\right) = \sqrt{\pi}$.
\end{itemize}

\paragraph*{Examples of Gamma Distribution}
\begin{itemize}
    \item Time until the $r$-th customer gets their coffee
    \item Milage that you get out of $r$ tanks of gas - note that $r$ does not have to be an integer here
\end{itemize}


\paragraph*{Relationship between Gamma and Exponential Distributions}
Much like the negative binomial gets us the probability of the $r$-th success in a series of Bernoulli trials, 
the Gamma distribution gives us the probability of the $r$-th event occurring in an exponential process.\\
\\
If $X_1, X_2, \ldots, X_r$ are independent Exponential random variables with parameter $\lambda$, then
\[
T = \sum_{i=1}^r X_i \sim \Gamma(r, \lambda).
\]

\paragraph*{Mean and Variance:} If $X \sim \Gamma(r, \lambda)$, then
\[\mu = E(X) = \frac{r}{\lambda}, \quad \sigma^2 = Var(X) = \frac{r}{\lambda^2}.\]

Note that the textbook defines the following:
\[ \alpha = r, \quad \beta = \frac{1}{\lambda}.\]

\paragraph*{The CDF of the Gamma Distribution}
For and arbitrary $r > 0$ there is no closed form for the CDF, but for positive integers $r$ we have
\[F(x) = P(X \leq x) = 1 - \sum_{k=0}^{r-1} \frac{(\lambda x)^k e^{-\lambda x}}{k!}.\]